\hypertarget{appendix-1dd4ad54ca57}{}
\section{Quadric error mesh simplification}\label{appendix-1dd4ad54ca57-quadric-error-mesh-simplification}

\textbf{Curriculum depth:} Developed · \textbf{First seen:} University courses · \textbf{Latest development:} University courses

\subsection{Definition and mechanism}\label{appendix-1dd4ad54ca57-definition-and-mechanism}

Quadric error simplification greedily collapses the edge with the least predicted geometric error. For each triangle plane with unit normal \(n\) and offset \(d\), form \(v=(n_x,n_y,n_z,d)^T\) and \(K=vv^T\). A homogeneous point \(U=(x,y,z,1)^T\) has squared point-to-plane error \(U^TKU\). Sum the matrices of incident planes at each vertex; an edge candidate combines endpoint matrices. Partitioning the combined quadratic gives \(Q(x)=x^TAx+2b^Tx+c\), whose stationary point solves \(Ax=-b\). If the system is singular or indefinite, test endpoints and midpoint and keep the least-cost candidate. After selecting a location, reject collapses that flip incident face normals. The metric is a local approximation, so topology preconditions and orientation checks remain necessary.

\subsection{Prerequisites}\label{appendix-1dd4ad54ca57-prerequisites}

\begin{itemize}
\tightlist
\item
  \href{../03-Visual-Computing-I/main.pdf\#nameddest=appendix-1aea26b2307f}{Half-edge mesh connectivity (03-Visual-Computing-I, appendix-1aea26b2307f)}
\item
  \hyperref[appendix-fab348609369]{Manifold polygon meshes and boundary structure}
\end{itemize}

\subsection{Used by}\label{appendix-1dd4ad54ca57-used-by}

\begin{itemize}
\tightlist
\item
  {Worlds, Images and Inference {[}referenced source unavailable{]}}
\end{itemize}

\subsection{Full teaching explanation}\label{appendix-1dd4ad54ca57-full-teaching-explanation}

\href{../03-Visual-Computing-I/main.pdf\#nameddest=03-visual-computing-i-ch002-meshes-connectivity-and-geometry-processing}{Meshes, connectivity, and geometry processing (03-Visual-Computing-I, 03-visual-computing-i-ch002-meshes-connectivity-and-geometry-processing)}
